\documentclass[11pt]{article}

\usepackage[margin=1in]{geometry}
\usepackage{parskip}
\usepackage{amsmath, amssymb, amsthm}
\usepackage{booktabs}
\usepackage{hyperref}

\title{A small library for the continued fraction expansion of rational
       numbers}
\author{Group N}
\date{}

\begin{document}
\maketitle

\begin{abstract}
We present \texttt{continued\_fractions}, a small Python library that
computes the continued fraction expansion of a rational number and the
corresponding list of convergents. The library is built on top of
\texttt{sympy}~\cite{meurer2017} so that every iterate is held exactly.
We illustrate the library on the rational approximations \(3/2\),
\(22/7\) and \(355/113\) of \(\pi\), and we recover the Fibonacci-ratio
approximations of the golden ratio. The documentation follows the
Diataxis framework~\cite{procida2017} and the conventions of the
open-source textbook \emph{Python for Mathematics}~\cite{knight2024}.
\end{abstract}

\section{Introduction}

The continued fraction expansion of a real number \(r\) is the (possibly
infinite) sequence of integers \(a_0, a_1, a_2, \dots\) such that

\[
    r = a_0 + \cfrac{1}{a_1 + \cfrac{1}{a_2 + \cfrac{1}{a_3 + \cdots}}}.
\]

For a rational number the expansion terminates, while for an irrational
number it is infinite. The convergents are the rational numbers obtained
by truncating the expansion after each step; they are known to give the
best rational approximations to \(r\) in the sense that no fraction with a
smaller denominator is closer~\cite{khinchin1964}.

We wanted a small library that we could use to inspect these expansions
without losing precision to floating-point rounding. The standard
library does not provide one. We chose to build on top of \texttt{sympy}
so that the iterates stay symbolic; this is the same approach taken by
the textbook \emph{Python for Mathematics}~\cite{knight2024}.

We present a three-function library,
\texttt{continued\_fractions}; we test it on the well-known rational
approximations \(22/7\) and \(355/113\) of \(\pi\); and we recover the
ratios of consecutive Fibonacci numbers as the convergents of the
golden ratio.

\section{The library}

The library lives in the single module \texttt{continued\_fractions.py}
and exposes three functions.

\texttt{continued\_fraction(rational\_number,
maximum\_terms)} computes the expansion of a rational by the Euclidean
algorithm: at each step we take the integer part of what is left,
subtract it, and invert the fractional part. The loop terminates as
soon as the remainder is zero, so we never run past the natural end of
the expansion.

\texttt{convergents(coefficients)} turns a
list of integer coefficients into the corresponding rationals. We use
the standard recursion

\[
    \frac{p_n}{q_n}
    = \frac{a_n p_{n - 1} + p_{n - 2}}{a_n q_{n - 1} + q_{n - 2}}.
\]

This avoids rebuilding the nested fraction from scratch at each step.

\texttt{golden\_ratio\_convergents(number\_of\_terms)} returns the
first convergents of the expansion \([1; 1, 1, 1, \dots]\). We use this
as a worked example because the convergents are exactly the ratios of
consecutive Fibonacci numbers.

\section{Worked examples}

Two common rational
approximations of \(\pi\) are \(22/7\) and \(355/113\). Their
expansions are \([3; 7]\) and \([3; 7, 16]\) respectively, and the
convergents recover the original rationals exactly, as shown in
Table~\ref{tab:pi}.

\begin{table}[h]
\centering
\begin{tabular}{l l l}
\toprule
Number & Expansion & Convergents \\
\midrule
\(22/7\)    & \([3; 7]\)     & \(3, 22/7\) \\
\(355/113\) & \([3; 7, 16]\) & \(3, 22/7, 355/113\) \\
\bottomrule
\end{tabular}
\caption{\textbf{Two rational approximations of \(\pi\).} For each
rational we show the continued fraction expansion and the
corresponding list of convergents. The last convergent agrees with the
input rational exactly.}
\label{tab:pi}
\end{table}

The continued fraction
\([1; 1, 1, 1, \dots]\) gives the convergents \(1, 2, 3/2, 5/3, 8/5,
13/8, \dots\), which are ratios of consecutive Fibonacci numbers. The
twentieth convergent agrees with \(\varphi = (1 + \sqrt{5}) / 2\) to
within \(10^{-8}\).

\section{Discussion}

The Euclidean step inverts a fraction, which on
floating-point arithmetic introduces a small error at each iterate. The
error compounds; by the third or fourth step it is large enough to
affect the integer part. Working with \texttt{sympy.Rational} sidesteps
this entirely. The cost is that the iterates can become large symbolic
fractions, but for the inputs we care about (rationals with small
denominators) this is not an issue.

We only handle rational inputs. A natural
extension would be to handle the periodic continued fractions of
quadratic irrationals such as \(\sqrt{2} = [1; 2, 2, 2, \dots]\), which
have a known closed form. This is a direction for future work.

\section{Conclusion}

We have presented \texttt{continued\_fractions}, a small Python library
for the continued fraction expansion of rational numbers. The library
recovers the standard rational approximations of \(\pi\), and the
convergents of the golden ratio match the ratios of consecutive
Fibonacci numbers to within tolerance.

\bibliographystyle{plain}
\bibliography{references}

\end{document}
